%% ============================================================================
%%  Library-Independent Adjudication of Three Rigid-Body Dynamics Engines
%%  Draft v0.3 -- blanks filled from the run record.
%%  Compile: pdflatex paper_silent_failure.tex  (twice for refs/ToC)
%% ============================================================================

\documentclass[11pt,a4paper]{article}

\usepackage[T1]{fontenc}
\usepackage[utf8]{inputenc}
\usepackage{amsmath,amssymb}
\let\Bbbk\relax
\usepackage{newtxtext,newtxmath}
\usepackage[margin=1.05in]{geometry}
\usepackage{booktabs}
\usepackage{longtable}
\usepackage{array}
\usepackage[expansion=false]{microtype}
\usepackage{enumitem}
\usepackage{listings}
\usepackage{xcolor}
\usepackage{colortbl}
\usepackage{titlesec}
\usepackage{caption}
\usepackage[colorlinks=true,linkcolor=linknavy,urlcolor=linknavy,citecolor=linknavy]{hyperref}
\newcommand{\doi}[1]{\href{https://doi.org/#1}{doi:#1}}

\definecolor{linknavy}{RGB}{20,60,110}
\definecolor{codebg}{RGB}{247,247,244}
\definecolor{coderule}{RGB}{205,205,198}
\definecolor{warnrule}{RGB}{198,150,80}
\definecolor{todored}{RGB}{150,45,45}

\titleformat{\section}{\normalfont\large\bfseries}{\thesection}{0.7em}{}
\titleformat{\subsection}{\normalfont\normalsize\bfseries}{\thesubsection}{0.6em}{}
\titlespacing*{\section}{0pt}{2.0ex plus .4ex}{1.0ex}

\renewcommand{\arraystretch}{1.15}
\newcolumntype{L}[1]{>{\raggedright\arraybackslash}p{#1}}
\captionsetup{font=small,labelfont=bf}

\lstdefinestyle{sh}{basicstyle=\ttfamily\small,
  backgroundcolor=\color{codebg}, frame=single,
  rulecolor=\color{coderule}, framesep=6pt, breaklines=true,
  showstringspaces=false, columns=fullflexible}
\lstset{style=sh}

\newenvironment{callout}
  {\par\medskip\noindent\begin{tabular}{@{}p{0.94\textwidth}@{}}
   \arrayrulecolor{warnrule}\hline\rule{0pt}{2.6ex}}
  {\\[0.8ex]\hline\arrayrulecolor{black}\end{tabular}\par\medskip}

%% Author-facing marker for the one item that cannot be filled by measurement.
\newcommand{\authortodo}[1]{\textcolor{todored}{\textbf{[#1]}}}

\title{\bfseries Library-Independent Adjudication of Three Rigid-Body\\
Dynamics Engines: Agreement, Divergence in Reach,\\
and a Shared Failure at Unstable Equilibria}

\author{Noah Parsons\\[0.3em]
{\normalsize Independent Researcher, Newcastle, Wyoming}\\
{\normalsize ORCID \href{https://orcid.org/0009-0000-7224-6040}{0009-0000-7224-6040}}}

\date{Draft v0.3 --- \today}

\begin{document}
\maketitle

\begin{abstract}
\noindent
Software that derives and integrates equations of motion is widely trusted to be
correct, yet independent implementations are rarely compared against one another
under adjudication that shares no code with any of them. This paper reports a
differential study of three rigid-body dynamics engines --- a symbolic compiler,
a general-purpose computer-algebra mechanics package, and a recursive
articulated-body library --- evaluated against a closed-form reference
implemented in NumPy alone that shares no library with any engine under test.

A 55-case adversarial suite spanning six axes was frozen before measurement, and
the comparison was extended to three structurally distinct mechanisms: an
$N$-link chain exercised from its regular regime into chaos, a slider--crank
whose closed loop and prismatic joint produce dead-centre configurations where
the effective inertia varies by four orders of magnitude over one revolution,
and a cart--pole whose mass-matrix determinant collapses by six orders of
magnitude twice per swing. Integration was pinned to a single scheme with
identical tolerances for all engines, which removes the integrator as a variable
and, by construction, forecloses any ranking of the engines on integration
accuracy.

The engines agree with the reference and with one another to within
$3\times10^{-14}$ on equations of motion and indistinguishably on energy
conservation, across both dynamical regimes and across constrained and
unconstrained pathways. Where they differ, they differ in reach rather than in
correctness: the three return equations for chains of 5, 12 and 30 links
respectively, and every engine that returned an answer returned a correct one.
No engine was observed to return a wrong answer while reporting success.

One failure was observed, and it is shared. At an unstable equilibrium, all
three engines and the independent reference report success while returning
approximately 20 radians of motion for a configuration whose exact solution does
not move; none warns that the trajectory is dominated by round-off amplified by
the instability. Silent failure in these tools is therefore a property of
finite-precision dynamics rather than a property that distinguishes
implementations --- a result that constrains what differential testing between
implementations can be expected to reveal about correctness in this domain.
\end{abstract}

\tableofcontents
\bigskip

% ===========================================================================
\section{Introduction}
% ===========================================================================

\subsection{The problem}
Scientific results increasingly rest on software that derives the equations
governing a physical system and integrates them forward. A researcher modelling
a joint, a mechanism, or a robot selects an engine, states the system, and takes
the resulting trajectory as an account of what the physical system would do. The
correctness of that account is rarely tested directly. It is inferred from the
engine's reputation, from its test suite, and from the absence of an error
message.

\subsection{The failure mode that motivates the study}
An engine that raises an exception is not dangerous; the researcher knows to
stop. The dangerous case is the engine that returns a plausible trajectory that
is wrong, with no indication that anything has gone amiss. Call this a
\emph{silent failure}. A silent failure is by construction the case nobody
thought to check for, which is why it survives in test suites written by the
same people who wrote the implementation.

\subsection{Two claims, often conflated}\label{sec:claims}
Discussion of this failure mode tends to run together two distinct propositions:

\begin{itemize}[leftmargin=2.2em,itemsep=0.25em]
  \item[\textbf{(A)}] These engines can report success while returning wrong
        physics.
  \item[\textbf{(B)}] Engines \emph{differ} in whether they do so, such that
        comparing implementations reveals incorrectness.
\end{itemize}

\noindent Claim (B) is the premise of differential testing: if two independent
implementations disagree, at least one is wrong, and the disagreement is the
signal. This paper set out to test (B) and did not find it.

\subsection{What was done}
Three independently developed engines were compared against a closed-form
reference implemented in NumPy alone, sharing no library with any of them. A
55-case adversarial matrix was frozen before measurement under a governance rule
forbidding removal of cases after their results were seen. Three mechanism
families were examined, chosen so that each degenerates by a different
mechanism. Integration was pinned to one scheme with identical tolerances for
every engine. The author's conflict of interest is disclosed in
\S\ref{sec:disclosure} and mitigated structurally rather than by declaration.

\subsection{What was found}
Agreement to the limit of double-precision arithmetic on every adjudicated case
in every family. Divergence between the engines in \emph{reach} --- the largest
system each will answer --- but not in correctness. One failure, at an unstable
equilibrium, shared by all three engines and by the independent reference alike.

\subsection{Why a negative result here is informative}
The measurement was designed to detect disagreement and demonstrably can. It
separated the engines by reach at $N=8$ and $N=20$
(\S\ref{sec:reach}), and it detected \emph{four} convention mismatches within
the measuring apparatus itself before any of them reached the reported results
(\S\ref{sec:threats}, \S\ref{sec:diagnostic}). An instrument that catches four
faults in its own construction, and that resolves a factor-of-six capability
difference between engines, is not an instrument that would miss a correctness
disagreement of the size differential testing is meant to expose. The absence of
such a disagreement is therefore an observation about the engines rather than a
limitation of the instrument.

\subsection{Contributions}
\begin{enumerate}[leftmargin=1.6em,itemsep=0.3em]
  \item A differential comparison of three independently developed rigid-body
        dynamics engines under adjudication by a reference sharing no library
        with any of them, across three mechanism families and both dynamical
        regimes.
  \item A negative result: no correctness disagreement was found on any axis
        measured; the engines separate by reach, not by honesty.
  \item A shared blind spot at unstable equilibria, present in all three engines
        \emph{and} in the independent reference, showing that the failure mode
        is a property of finite-precision dynamics rather than of any
        implementation.
  \item A methodological account --- frozen case matrix, library-independent
        referee, integrator pinning, maintainer disclosure --- reusable for
        cross-implementation comparison in other computational domains.
  \item A transferable diagnostic for convention mismatch in
        cross-implementation comparison, derived from three independent
        instances encountered in this study.
\end{enumerate}

% ===========================================================================
\section{Related work}
% ===========================================================================

\subsection{Differential testing and the independence assumption}
Differential testing --- running two or more implementations of the same
specification on the same input and treating any difference in output as
evidence of a defect --- was named and systematised by
McKeeman~\cite{mckeeman1998}, and has since become a standard instrument
wherever a specification admits multiple implementations. Its appeal is that it
dissolves the test oracle problem: no statement of correct behaviour is needed,
only two implementations and a disagreement.

Its premise is that independent implementations fail independently. That premise
has been tested directly, and it failed. Knight and Leveson~\cite{knight1986}
prepared 27 versions of one program independently from a common specification at
two universities and subjected them to one million tests. The versions were
individually very reliable, but the number of tests on which more than one
version failed was substantially greater than independence predicts. Their
result is the standard caution against $N$-version programming, and it is the
closest antecedent of the present paper's central finding.

The mechanism differs, and the difference is the contribution. Knight and Leveson
attributed correlated failure to shared human misconception: independent
programmers reading the same specification make the same mistakes, because some
parts of a problem are harder to think about than others. The correlated failure
reported in \S7 has no human cause at all. Every implementation examined here ---
including a reference written from first principles specifically to be
independent --- fails identically at an unstable equilibrium, because the failure
originates in floating-point arithmetic rather than in anyone's reasoning. Shared
misconceptions can in principle be diluted by adding more independent authors.
Shared arithmetic cannot.

\subsection{The oracle problem and scientific software}
Barr et al.~\cite{barr2015} survey the test oracle problem and classify the
partial remedies available when no exact oracle exists: derived oracles,
metamorphic relations, and \emph{pseudo-oracles} --- independently produced
implementations used as a standard of comparison. The reference used here is a
pseudo-oracle in that sense, distinguished by an explicit and checkable
independence criterion: it shares no library with any engine under test, which is
a stronger condition than independent authorship and a much stronger one than
independent implementation within a common ecosystem.

Kanewala and Bieman~\cite{kanewala2014} review testing practice in scientific
software specifically and identify the oracle problem as one of its two defining
difficulties: for software whose purpose is to compute a result nobody knows in
advance, the expected output is frequently unavailable by construction. This is
precisely the situation of a dynamics engine asked for the trajectory of a
mechanism that has never been built.

The most direct empirical precedent is Hatton and Roberts~\cite{hatton1994}, who
ran nine independently developed seismic data processing packages on identical
inputs and found agreement to roughly one significant figure --- a disagreement
so large that it called into question published results in the field. Their study
and this one share a design and reach opposite conclusions, which requires
explanation rather than assertion. Three differences are plausible. Their
packages implemented a long processing pipeline whose specification left
considerable latitude, whereas the equations of motion of a rigid mechanism are
mathematically determined once the Lagrangian is fixed. Their comparison was
engine-against-engine, with no external standard, so disagreement could be
localised but not attributed. And three decades separate the two studies. The
present result should therefore be read narrowly: rigid-body dynamics is a domain
where the specification is unusually tight, and the finding of agreement does not
transfer to computational domains where it is not.

\subsection{Verification and validation of multibody codes}
The multibody dynamics community maintains its own verification infrastructure.
The IFToMM Technical Committee for Multibody Dynamics hosts a library of
computational benchmark problems~\cite{iftomm} intended to standardise
comparison of simulation software on both accuracy and efficiency. That effort is
complementary to this one but differs in its notion of ground truth: benchmark
solutions are established by consensus among submitted solvers, which is a
reasonable practice and a different epistemic basis from a closed-form reference
that shares no code with any participant.

Nearest in subject matter, Gonz\'alez et al.~\cite{gonzalez2016} examine the
behaviour of augmented-Lagrangian and Hamiltonian formulations for constrained
multibody systems in the neighbourhood of singular configurations, identifying
the sources of numerical difficulty there and comparing formulations for
robustness. Their singular configurations are the analogue of the dead centres of
\S4.2. The distinction is that they compare \emph{formulations} within a
controlled implementation, isolating the effect of mathematical choice, whereas
this study compares \emph{implementations} of nominally the same formulation,
isolating the effect of engineering. The two questions are separable and both
worth asking; a formulation that is robust in principle can still be implemented
incorrectly, and this paper measures the second.

\subsection{Comparisons of dynamics engines}
Erez, Tassa and Todorov~\cite{erez2015} compare five simulation engines --- among
them Bullet, MuJoCo, ODE and PhysX --- on quantitative measures chosen for
robotics workloads, instantiating the same model in each and running side-by-side.
Their principal conclusion is that each engine performs best on the class of
system it was designed and optimised for. That finding is reproduced
incidentally here, and by a different measurement: the recursive
articulated-body library returns the largest residuals of the three engines on
long serial chains and the smallest on the two-body tree
(\S\ref{sec:cartpole-results}), an ordering that reverses with the shape of the
problem rather than with the quality of the engine.

The difference in aim is what leaves room for the present study. Erez et al.\
measure speed, stability, and behaviour under contact --- properties that matter
for whether an engine is usable --- and adjudicate against no external standard,
because for the contact-rich systems they consider none exists in closed form.
This paper asks a narrower question, on systems chosen precisely because closed
forms do exist: not whether the engines are usable but whether their equations
are right, judged against a derivation that shares nothing with any of them.

\subsection{The gap}
Prior work establishes that differential testing is powerful, that its
independence premise is unreliable in the presence of shared human error, that
scientific software has historically disagreed badly under $N$-version
comparison, and that dynamics engines differ measurably in performance. What has
not been reported, to the author's knowledge, is a comparison of symbolic and
numerical dynamics engines adjudicated by a closed-form reference sharing no
library with any of them, over a case matrix frozen before measurement, with the
resulting disagreements characterised. That is the gap this paper addresses, and
the result --- agreement everywhere except in a failure the referee shares --- is
informative chiefly because the design would have detected the alternative.

% ===========================================================================
\section{Method}
% ===========================================================================

\subsection{The engines under test}
Three engines were selected on the criterion that each derives the equations of
motion of a mechanical system by a different route, and that no two share a
substantial body of code.

\begin{center}
\begin{tabular}{@{}lL{0.30\textwidth}L{0.26\textwidth}@{}}
\toprule
\textbf{Engine} & \textbf{Route} & \textbf{Pathways exercised} \\
\midrule
MechanicsDSL 2.1.3 (\texttt{a8dc2b2}) & symbolic compiler; consumes a Lagrangian, direct solve for a fixed topology & Lagrangian, Hamiltonian, constrained \\
\texttt{sympy.physics.mechanics} 1.14.0 & general-purpose computer algebra, driven as documented & Lagrangian, constrained \\
Drake 1.56.0 & recursive articulated-body library for arbitrary topologies, contact, real-time control & model-building interface \\
\bottomrule
\end{tabular}
\end{center}

Versions were pinned for the duration of the study. Pinning is not a formality:
an engine that changes between the derivation and integration phases of a
comparison invalidates the comparison, and the version under test must be
recoverable by a reader attempting to reproduce the measurement.

Each engine was exercised through the pathways it exposes. Two support a
Lagrangian formulation; one additionally supports a Hamiltonian formulation; one
constructs models through a model-building interface rather than from a
Lagrangian. This asymmetry is a property of the engines and is reported rather
than normalised away, since normalising it would require driving at least one
engine in a manner its documentation does not describe.

\subsection{The referee}
Adjudication is performed by a closed-form reference derived by hand and
implemented in NumPy alone. The reference shares no library with any engine
under test. This is the central design decision of the study, and its rationale
is straightforward: two implementations that share a symbolic algebra layer, a
linear solver, or a code-generation backend may agree because they inherit the
same defect, and their agreement then carries no information about correctness.
A referee that shares nothing with the engines cannot agree with them for that
reason.

For the unconstrained families the reference solves the equations of motion
directly for a fixed topology. For the constrained families both systems admit a
constant diagonal mass matrix and holonomic constraints, so accelerations and
multipliers follow from the saddle-point system
\begin{equation}
\begin{bmatrix} M & J^{\top} \\ J & 0 \end{bmatrix}
\begin{bmatrix} \ddot q \\ \lambda \end{bmatrix}
=
\begin{bmatrix} F(q) \\ -\dot J \dot q \end{bmatrix},
\label{eq:kkt}
\end{equation}
in which both $J$ and $\dot J\dot q$ have closed forms for the rod and circle
constraints involved. No symbolic algebra is required at any point in the
reference, which is what permits the claim of library independence to be made
literally rather than approximately.

The reference is deliberately left unstabilised: no Baumgarte term and no
projection. A stabilised reference would be solving different equations from
those the engines are asked for, and the resulting comparison would be between
two different problems. The cost of that choice is a constraint drift of order
$10^{-4}$ over three seconds, which does not affect the derivation comparison
because that comparison evaluates accelerations at given states.

\subsection{The case matrix and the freeze}
The case matrix comprises 55 cases over 36 systems across six axes: degrees of
freedom, closed-loop topology, constraint redundancy, near-singular inertia,
mass ratio, and symbolic nesting depth. The matrix was frozen --- fixed in
content and recorded with a version identifier --- before the cross-engine
measurements reported here were taken.

The freeze is governed by an explicit rule. Extension of the harness is
permitted where it adds adjudication without revising existing verdicts, and any
new axis must be declared before it is measured. Removing a case after observing
its result is selection on the outcome and is forbidden. The rule exists because
the author maintains one of the engines under test, and the temptation it guards
against is therefore not hypothetical.

Of the 55 cases, 45 are adjudicated; the remainder exceed the wall-clock limit
for at least one engine and are recorded as timeouts rather than as verdicts. A
timeout is scored as a distinct category from a pass or a failure, because an
absence of an answer and a wrong answer are different events and conflating them
would obscure precisely the distinction this study exists to examine.

\subsection{Integrator pinning}\label{sec:pinning}
Every engine supplies only its right-hand side. Integration is performed by a
single call to one explicit Runge--Kutta scheme (DOP853) with identical relative
and absolute tolerances ($\mathrm{rtol}=10^{-10}$, $\mathrm{atol}=10^{-12}$) for
all engines.

Pinning has a consequence that must be stated plainly: it makes the measurement
structurally incapable of ranking the engines on integration accuracy. Energy
drift under a pinned integrator is dominated by the integrator the engines
share, not by the engines. This is not a limitation of the design but its
purpose. Identical drift across engines is evidence that their equations are
equivalent, which is what the comparison exists to establish. Without pinning,
one engine's native integrator would have been compared against another's, and
the difference attributed to the engines.

The design has a corresponding limitation, disclosed here rather than in the
results. An engine whose constrained dynamics are enforced inside its own solver
cannot participate in a pinned-integrator comparison on constrained systems
without being asked to do something other than what it is built to do. Where
this arises, the affected engine is scored on constraint satisfaction alone, and
the asymmetry is reported rather than concealed. This is the case for the
recursive library on the slider--crank (\S\ref{sec:slidercrank-results}); the
same engine participates fully on the cart--pole, which is a tree and requires
no constraint.

\subsection{Adjudication criteria}
Two quantities are compared. The first is the residual between an engine's
computed accelerations and the reference's, evaluated at 16 probe states per
case drawn from the constraint manifold where one applies. The second is energy
drift over a fixed integration horizon under the pinned integrator. For
constrained cases a third quantity, constraint residual drift from $t=0$, is
recorded. The adjudication tolerance throughout is $10^{-8}$ relative.

Constraint residual is measured as drift from the initial value rather than as
an absolute quantity. The suite's constrained initial conditions are rounded to
four decimal places, so every constrained case begins approximately $10^{-4}$
off its own constraint manifold. This is harmless for the classifier as
constructed, but a reader checking the absolute residual would find an apparent
violation and might reasonably attribute it to the engine. The fact is recorded
here so that it cannot mislead.

\subsection{Threats to validity}\label{sec:threats}

\paragraph{Common authorship of engine and referee.}
The referee shares no library with any engine, but it does share an author with
one of them. A misconception held by that author could in principle appear in
both, and their agreement would then be uninformative. Two considerations bound
this risk. The two engines the author did not write agree with the reference to
the same tolerance, and a shared misconception would have to be present in all
four independently.

Separately, the study record documents \textbf{four} instances in which the
author's expectation about the \emph{measuring apparatus} was contradicted by an
independent artefact, and corrected before the affected results were reported:

\begin{center}
\begin{tabular}{@{}lL{0.42\textwidth}L{0.22\textwidth}@{}}
\toprule
\textbf{\#} & \textbf{Mistaken expectation} & \textbf{Caught by} \\
\midrule
1 & The Hamiltonian pathway's right-hand side returns accelerations; it returns $\dot p$ & discriminating $\dot q$ against $M^{-1}p$ \\
2 & A revolute chain reports absolute joint angles; it reports relative ones & inspecting the mass matrix \\
3 & A pole offset rotated about $+y$ matches a reference using $x + l\sin\theta$; it mirrors it & inspecting the mass matrix \\
4 & An engine reporting motion at an exact equilibrium is wrong; the exact solution is unreachable in finite precision & the independent reference behaving identically \\
\bottomrule
\end{tabular}
\end{center}

\noindent Instances 1--3 are convention mismatches in the comparison and are the
basis of the diagnostic in \S\ref{sec:diagnostic}; instance 4 is a mistaken
adjudication criterion. All four concern the instrument, and all four were
detected by consulting an artefact rather than by reasoning further.

The study record contains further corrected errors that concern external facts
rather than the instrument --- chiefly the provenance of a third-party API
guard, established only after documentation, a cached search result, and an
unvalidated repository query had each produced a different wrong answer. Those
are omitted here. They bear on the author's diligence, not on whether the
measurement detects its own faults, and only the latter is evidence.

\paragraph{Mechanism coverage.}
Three mechanism families are examined. They were chosen for structurally
distinct degeneracies --- a vanishing transmission ratio, a collapsing
mass-matrix determinant, and a chaotic regime reached by amplitude --- but they
remain three families, and the results do not license claims about mechanism
classes outside them.

\paragraph{A comparator that is not independent on one family.}
On the cart--pole the computer-algebra package's residual against the reference
is \emph{exactly zero at all twenty configurations}, not merely small. This is
addressed in \S\ref{sec:cartpole-results}: on that family the two are not
independent, and the effective comparison is three-way rather than four-way.

\paragraph{Wall-clock timeouts.}
The reach results in \S\ref{sec:reach} rest on a wall-clock limit of 120\,s
measured on a single machine in single runs. The ordering of the engines is
robust to plausible variation in that limit, since the separations are of order
minutes rather than seconds, but the specific link counts at which each engine
ceases to return are machine-dependent and should be read as such.

\paragraph{Non-portable axes.}
One axis of the frozen matrix --- symbolic nesting depth --- is meaningful only
to an engine that consumes raw Lagrangians, and no library-independent reference
for it can exist. Those cases are excluded from the cross-engine claim by
construction rather than by choice.

\paragraph{Constrained continuous-time dynamics in the recursive library.}
Constrained dynamics in the recursive library are enforced within its own
discrete solver and are not available through the continuous-time interface used
elsewhere in this study. The engine is therefore driven in its discrete mode on
the slider--crank, where it cannot share the pinned integrator and is scored on
constraint satisfaction alone. This is a property of the engine's supported
interfaces and is reported as a limitation of the comparison; no defect claim is
made here.

\subsection{Disclosure}\label{sec:disclosure-method}
The author maintains one of the three engines measured. The mitigation is
structural rather than declarative: no engine adjudicates itself or any other,
the referee shares no library with any of them, the case matrix was frozen
before measurement under a rule forbidding post hoc removal, and the results
reported include the measurements on which the author's engine performs worst
(\S\ref{sec:behind}).

AI assistance was used for portions of the implementation and analysis; see the
disclosure statement.

All figures reported were produced by executing committed code. Artefacts ---
the case matrix, the reference implementation, the adapters, the sweep scripts,
and the raw output records --- are archived at
\href{https://doi.org/10.5281/zenodo.22315173}{\texttt{doi:10.5281/zenodo.22315173}}
under the MIT licence.

% ===========================================================================
\section{Systems}
% ===========================================================================

Three mechanism families are examined. They were chosen so that each presents a
degeneracy arising from a different mechanism, on the reasoning that a
comparison exercising one kind of hard case repeatedly establishes less than one
exercising several kinds once.

\subsection{The planar \texorpdfstring{$N$}{N}-link chain}
An open chain of $N$ links under gravity, with revolute joints and no
constraint. The chain scales in degrees of freedom without changing its
structure, which makes it the natural vehicle for the reach measurements of
\S\ref{sec:reach}, and its dynamics pass from regular to chaotic as the initial
amplitude increases, which makes it the vehicle for the amplitude axis.

The degeneracy here is dynamical rather than geometric. At the suite's original
initial angles the motion is quasi-periodic: perturbing the initial state by
$10^{-10}$ and integrating for ten seconds produces growth of order $10^{0}$. At
larger amplitudes the same perturbation grows by eleven orders of magnitude over
the same horizon. Growth of order $10^{0}$ is polynomial and indicates regular
motion; growth of order $10^{11}$ is exponential and indicates chaos.

The amplitude axis was added after measurement showed that the frozen suite had
been exercising the chain only in its regular regime --- that is, that a suite
whose stated method is to increase difficulty until behaviour breaks down had
been operating entirely within the well-behaved region of its principal system.
The axis was declared before it was measured, in accordance with the governance
rule, and the Lagrangian it uses is character-for-character the one already
frozen; only the initial conditions differ.

\subsection{The slider--crank}
A crank rotating about a fixed pivot, connected by a rigid rod to a slider
constrained to a straight line. The mechanism introduces two features absent
from the chain: a closed kinematic loop, and a prismatic joint alongside a
revolute one.

The constraint admits a closed form. Expanding $|S-P|^2 - l^2$ yields
$x^2 - 2rx\cos\theta + r^2 - l^2$, the trigonometric identity eliminating the
remaining terms, so no symbolic algebra is required to construct the reference.

The degeneracy is geometric and arises from the transmission ratio. The
effective inertia
\begin{equation}
M_{\mathrm{eff}}(\theta) = m_c r^2 + m_s\left(\frac{\mathrm{d}x}{\mathrm{d}\theta}\right)^2
\label{eq:meff}
\end{equation}
varies by four orders of magnitude over a single revolution when the slider is
heavy, because $\mathrm{d}x/\mathrm{d}\theta$ vanishes exactly at $\theta=0$ and
$\theta=\pi$. At those two configurations --- the dead centres --- the slider
ceases to contribute inertia altogether and the mechanism approaches the loss of
its mass matrix. The measured variation is $1.12\times10^{2}$ at
$m_s/m_c = 10^{2}$, $1.11\times10^{4}$ at $10^{4}$, and $1.11\times10^{6}$ at
$10^{6}$. The degeneracy is produced by the geometry rather than by a
hand-constructed ill-conditioned matrix, which is what makes it a fair test.

Two parameters were swept: the rod-to-crank length ratio, from 3.0 down to 1.01
as the mechanism approaches its folding configuration, and the slider-to-crank
mass ratio to $10^{6}$. Engines were evaluated both at random states on the
constraint manifold and exactly at both dead centres.

\subsection{The cart--pole}
A cart translating on a horizontal prismatic joint with a pole suspended from it
by a revolute joint. Unlike the slider--crank the system is a tree, with no loop
and no constraint, which permits all three engines to be compared in their
ordinary modes of operation under a single pinned integrator, with nothing
excluded for reasons unrelated to dynamics.

The system is structurally new to the study in two respects: it combines joint
types within one mechanism, and its mass matrix is genuinely configuration
dependent, where the suite's earlier multi-body systems had constant ones.

Its degeneracy arises from a third mechanism again. The determinant
\begin{equation}
\det M(q) = m l^{2}\left(M + m\sin^{2}\theta\right)
\label{eq:detM}
\end{equation}
collapses by six orders of magnitude twice per swing when the cart is light and
the pole passes through vertical, because a near-massless cart cannot resist the
pole's horizontal reaction. The slider--crank degenerates because a transmission
ratio vanishes; the cart--pole degenerates because an inertia vanishes. The
causes are unrelated, which is the reason for testing both.

% ===========================================================================
\section{Results: agreement}
% ===========================================================================

\subsection{Equations of motion}
Residuals between each engine's computed accelerations and the reference's are
constant across the amplitude axis, as expected: the equations of motion do not
depend on where the system starts.

\begin{table}[h]
\centering
\caption{Equation-of-motion residuals against the reference, three-link chain,
16 probe states drawn uniformly from $[-0.5,0.5]$. Relative to the
double-precision unit round-off $u \approx 2.22\times10^{-16}$.}
\label{tab:eom}
\begin{tabular}{@{}lrr@{}}
\toprule
\textbf{Engine} & \textbf{Residual} & \textbf{Multiple of $u$} \\
\midrule
MechanicsDSL & $1.887\times10^{-15}$ & $8.5$ \\
Drake        & $1.072\times10^{-14}$ & $48$ \\
SymPy        & $2.620\times10^{-14}$ & $118$ \\
\bottomrule
\end{tabular}
\end{table}

All three engines are therefore correct to the limit of the arithmetic, and the
differences between them reflect accumulated rounding over differing operation
counts rather than any difference in the equations derived.

\subsection{The amplitude axis}
Across eight amplitudes spanning perturbation growth from single figures to
$5.4\times10^{11}$, and three engines, 24 engine--amplitude pairs were
evaluated. There were no refusals, no incorrect equations, and no integration
failures.

\begin{table}[h]
\centering
\caption{Amplitude axis, three-link chain. Perturbation growth is the factor by
which a $10^{-10}$ initial displacement grows over 10\,s, measured on the
reference at $\mathrm{rtol}=10^{-12}$. Energy drift is relative, under the
pinned integrator.}
\label{tab:amplitude}
\begin{tabular}{@{}rrlrrr@{}}
\toprule
\textbf{Amp.\ (rad)} & \textbf{Growth} & \textbf{Regime} & \textbf{MechanicsDSL} & \textbf{SymPy} & \textbf{Drake} \\
\midrule
$0.15$   & $2.3$                & regular & $7.302\times10^{-11}$ & $7.328\times10^{-11}$ & $7.445\times10^{-11}$ \\
$0.50$   & $3.9$                & regular & $2.201\times10^{-11}$ & $2.201\times10^{-11}$ & $2.201\times10^{-11}$ \\
$1.00$   & $3.8$                & regular & $2.702\times10^{-11}$ & $2.702\times10^{-11}$ & $5.123\times10^{-11}$ \\
$1.50$   & $3.5\times10^{4}$    & chaotic & $1.723\times10^{-10}$ & $1.723\times10^{-10}$ & $2.860\times10^{-10}$ \\
$2.00$   & $1.1\times10^{7}$    & chaotic & $2.137\times10^{-10}$ & $2.137\times10^{-10}$ & $2.139\times10^{-10}$ \\
$2.50$   & $2.8\times10^{10}$   & chaotic & $4.737\times10^{-10}$ & $4.767\times10^{-10}$ & $7.783\times10^{-10}$ \\
$3.00$   & $1.7\times10^{11}$   & chaotic & $3.036\times10^{-10}$ & $3.036\times10^{-10}$ & $3.042\times10^{-10}$ \\
$\pi$    & $5.4\times10^{11}$   & chaotic & $2.737\times10^{-10}$ & $2.737\times10^{-10}$ & $2.737\times10^{-10}$ \\
\bottomrule
\end{tabular}
\end{table}

Energy drift rises with amplitude, as harder dynamics demand more of the
integrator, and does so almost identically for all three engines: at the worst
case the three values are $4.74$, $4.77$ and $7.78\times10^{-10}$, the same
quantity to within a factor of $1.6$.

\subsection{The chaotic regime}\label{sec:chaotic}
At three degrees of freedom with all links started well into the chaotic regime,
energy drift is indistinguishable across the reference and all three engines.

\begin{table}[h]
\centering
\caption{Chaotic regime, three-link chain, all links started at $3.0$\,rad.
Perturbation growth $1.26\times10^{11}$ over 10\,s at
$\mathrm{rtol}=10^{-10}$, implying a Lyapunov rate
$\lambda = 2.556\,\mathrm{s}^{-1}$. Divergence is the time at which the
trajectory departs the reference's by $10^{-6}$; \S\ref{sec:divergence}
explains why it is not a ranking.}
\label{tab:chaotic}
\begin{tabular}{@{}lrr@{}}
\toprule
\textbf{Engine} & \textbf{Energy drift} & \textbf{Divergence} \\
\midrule
reference    & $3.0358\times10^{-10}$ & --- \\
MechanicsDSL & $3.0356\times10^{-10}$ & $6.96$\,s \\
SymPy        & $3.0361\times10^{-10}$ & $6.28$\,s \\
Drake        & $3.0421\times10^{-10}$ & $5.69$\,s \\
\bottomrule
\end{tabular}
\end{table}

\noindent The largest separation is $0.2$ per cent relative, twelve orders of
magnitude below any physically meaningful quantity.

\subsection{The constrained pathway}\label{sec:constrained-results}
All eight constrained cases in the adjudicated set agree with the closed-form
reference of equation~\eqref{eq:kkt}.

\begin{table}[h]
\centering
\caption{Constrained pathway, MechanicsDSL against the saddle-point reference,
16 probe states per case. \texttt{redund\_R$k$} carries $k+1$ constraints on two
coordinates, all scalar multiples of one another.}
\label{tab:constrained}
\begin{tabular}{@{}lrrr@{}}
\toprule
\textbf{Case} & \textbf{Coordinates} & \textbf{Constraints} & \textbf{Residual} \\
\midrule
\texttt{loops\_N3}  & 4 & 3 & $5.788\times10^{-16}$ \\
\texttt{redund\_R0} & 2 & 1 & $1.958\times10^{-16}$ \\
\texttt{redund\_R1} & 2 & 2 & $1.084\times10^{-15}$ \\
\texttt{redund\_R2} & 2 & 3 & $1.234\times10^{-14}$ \\
\texttt{redund\_R3} & 2 & 4 & $2.554\times10^{-15}$ \\
\texttt{redund\_R4} & 2 & 5 & $2.004\times10^{-15}$ \\
\texttt{redund\_R6} & 2 & 7 & $3.266\times10^{-15}$ \\
\texttt{redund\_R8} & 2 & 9 & $7.230\times10^{-15}$ \\
\bottomrule
\end{tabular}
\end{table}

The redundancy family is the more informative half. Its constraint Jacobian
retains rank one while its row count grows to nine, so the multipliers become
underdetermined while the accelerations remain unique. The engines track the
reference throughout. An exactly redundant constraint set does not determine how
the constraint force is distributed, only its total, and neither the engines nor
the reference is troubled by this.

\paragraph{Coverage.}
With these cases adjudicated, the number of adjudicated cases carrying an
independent derivation check rises from \textbf{22 to 40} of 45. The arithmetic
does not compose as a simple increment, and the reconciliation is given in
Table~\ref{tab:coverage} because the figure is citable.

\begin{table}[h]
\centering
\caption{Coverage reconciliation over the 45 adjudicated cases. The two referees
cover overlapping but different sets, so the increment is not additive.}
\label{tab:coverage}
\begin{tabular}{@{}lrL{0.40\textwidth}@{}}
\toprule
\textbf{Referee} & \textbf{Cases} & \textbf{Composition} \\
\midrule
Computer-algebra oracle (original) & 22 & Lagrangian pathway only, including 5 symbolic-axis cases \\
Library-independent reference      & 35 & \texttt{dof} 7, \texttt{mass\_ratio} 14, \texttt{redundancy} 7, \texttt{near\_singular} 6, \texttt{loops} 1 \\
\midrule
\textbf{Union: any independent check} & \textbf{40} & \\
\quad gained by the reference & 18 & \\
\quad reachable only by the original oracle & 5 & symbolic-axis Lagrangian cases \\
\midrule
\textbf{Carrying no independent check} & \textbf{5} & symbolic-axis Hamiltonian cases \\
\bottomrule
\end{tabular}
\end{table}

\noindent The five cases that remain unchecked are the Hamiltonian half of the
symbolic axis. No library-independent reference for that axis can exist: nested
cosine depth is meaningful only to an engine that consumes a raw Lagrangian, so
there is no system to derive independently.

\subsection{The slider--crank and the cart--pole}\label{sec:cartpole-results}
Across both additional families, 72 engine--case pairs were evaluated with no
disagreements and no refusals.

\begin{table}[h]
\centering
\caption{Cart--pole, five cart-to-pole mass ratios $\times$ four initial angles
$\times$ three engines. Residuals are constant across the four angles at each
mass ratio. Energy drift is the worst over the four angles.}
\label{tab:cartpole}
\begin{tabular}{@{}rrrrrr@{}}
\toprule
& & \multicolumn{2}{c}{\textbf{Residual}} & \multicolumn{2}{c}{} \\
\cmidrule(lr){3-4}
$M/m$ & $\det M$ range & MechanicsDSL & Drake & SymPy & Worst drift \\
\midrule
$10^{1}$  & $1.1$      & $3.72\times10^{-16}$ & $2.78\times10^{-17}$ & $0$ & $1.99\times10^{-11}$ \\
$10^{0}$  & $2.0$      & $3.72\times10^{-16}$ & $1.86\times10^{-16}$ & $0$ & $3.75\times10^{-10}$ \\
$10^{-2}$ & $10^{2}$   & $5.02\times10^{-15}$ & $2.05\times10^{-16}$ & $0$ & $1.60\times10^{-9}$ \\
$10^{-4}$ & $10^{4}$   & $7.49\times10^{-15}$ & $5.55\times10^{-17}$ & $0$ & $1.11\times10^{-9}$ \\
$10^{-6}$ & $10^{6}$   & $1.25\times10^{-14}$ & $1.96\times10^{-16}$ & $0$ & $3.36\times10^{-9}$ \\
\bottomrule
\end{tabular}
\end{table}

All three engines agree with the reference at every mass ratio and every
starting angle, including where the determinant collapses by six orders of
magnitude, and energy drift is of order $10^{-9}$ or better throughout.

\begin{callout}
\textbf{The computer-algebra package is not an independent comparator on this
family.} Its residual against the reference is \emph{exactly} zero --- not small
--- at all twenty configurations, over all sixteen probe states each. A residual
of exactly zero under a relative-error metric means the two implementations
produced bit-identical floating-point results, which indicates that they
executed the same arithmetic in the same order rather than that they agree by
independent derivation. The likely cause is that the cart--pole's mass matrix
and forcing vector are simple enough that the package's symbolic simplification
reproduces the hand-derived closed form exactly, and its code generation then
emits the same sequence of operations.

Consequently the cart--pole comparison is effectively three-way --- reference,
symbolic compiler, recursive library --- not four-way. This does not affect the
other two families, where the same package returns non-zero residuals
($2.62\times10^{-14}$ on the chain, $5.55\times10^{-16}$ to
$1.96\times10^{-15}$ on the slider--crank) and is a genuine independent
comparator.
\end{callout}

\noindent Subject to that qualification, the recursive library returns the
smallest residuals on this family, of order $10^{-16}$ --- better than the
symbolic compiler by two orders of magnitude. This is the expected outcome when
an algorithm designed for arbitrary topologies is applied to a two-body tree,
and it is the converse of the pattern observed on the longer chains, where its
residuals are the largest.

\begin{table}[h]
\centering
\caption{Slider--crank. Residual and energy drift under the pinned integrator
for the two engines that can share it; the recursive library is driven in its
discrete mode and scored on constraint satisfaction alone
(\S\ref{sec:pinning}). Constraint column is the worst rod-length residual over
a 5\,s integration.}
\label{tab:slidercrank}
\begin{tabular}{@{}rrlrrr@{}}
\toprule
$l/r$ & $m_s/m_c$ & \textbf{Engine} & \textbf{Residual} & \textbf{Energy drift} & \textbf{Constraint} \\
\midrule
$3$    & $10^{2}$ & MechanicsDSL   & $1.22\times10^{-15}$ & $1.51\times10^{-10}$ & $1.81\times10^{-10}$ \\
$3$    & $10^{2}$ & SymPy          & $5.55\times10^{-16}$ & $1.51\times10^{-10}$ & $1.81\times10^{-10}$ \\
$3$    & $10^{2}$ & Drake (disc.)  & ---                  & ---                  & $1.77\times10^{-6}$ \\
$1.5$  & $10^{2}$ & MechanicsDSL   & $4.17\times10^{-15}$ & $2.11\times10^{-10}$ & $9.12\times10^{-10}$ \\
$1.5$  & $10^{2}$ & SymPy          & $7.53\times10^{-16}$ & $2.12\times10^{-10}$ & $9.12\times10^{-10}$ \\
$1.5$  & $10^{2}$ & Drake (disc.)  & ---                  & ---                  & $1.25\times10^{-7}$ \\
$1.05$ & $10^{2}$ & MechanicsDSL   & $5.66\times10^{-15}$ & $5.40\times10^{-10}$ & $1.40\times10^{-9}$ \\
$1.05$ & $10^{2}$ & SymPy          & $1.96\times10^{-15}$ & $5.40\times10^{-10}$ & $1.40\times10^{-9}$ \\
$1.05$ & $10^{2}$ & Drake (disc.)  & ---                  & ---                  & $1.42\times10^{-7}$ \\
$3$    & $10^{4}$ & MechanicsDSL   & $1.97\times10^{-15}$ & $2.92\times10^{-10}$ & $8.15\times10^{-10}$ \\
$3$    & $10^{4}$ & SymPy          & $5.55\times10^{-16}$ & $2.90\times10^{-10}$ & $8.11\times10^{-10}$ \\
$3$    & $10^{4}$ & Drake (disc.)  & ---                  & ---                  & $1.14\times10^{-6}$ \\
\bottomrule
\end{tabular}
\end{table}

\label{sec:slidercrank-results}
Both symbolic engines hold the rod length to approximately $10^{-9}$ through
dead-centre crossings at a length ratio of $1.05$ and a mass ratio of $10^{4}$.
The recursive library holds the rod to approximately $10^{-6}$ in its discrete
mode.

\subsection{Summary}
No engine returned a wrong answer on any adjudicated case, in any family, in
either dynamical regime, on either the constrained or the unconstrained pathway.
Claim (B) of \S\ref{sec:claims} --- that engines differ in whether they report
success on incorrect physics --- is unsupported by every measurement in this
study.

% ===========================================================================
\section{Results: divergence in reach}\label{sec:reach}
% ===========================================================================

Agreement on cases all three engines can answer leaves open whether they can
answer the same cases at all. Each engine was therefore taken up a ladder of
system sizes to 30 links, with each engine--pathway--size combination built in
its own subprocess under a fixed 120\,s wall-clock limit, and the time measured
to produce the equations of motion and evaluate the accelerations once. Results
were checked against the reference wherever an answer was returned.

\begin{table}[h]
\centering
\caption{Scaling ladder, Lagrangian pathway, 120\,s wall clock per
engine--size pair. Single runs on one machine. Times include model construction
and one acceleration evaluation. Once an engine fails to return it is not
retried at larger sizes.}
\label{tab:scaling}
\begin{tabular}{@{}rlrlrlr@{}}
\toprule
& \multicolumn{2}{c}{\textbf{MechanicsDSL}} & \multicolumn{2}{c}{\textbf{SymPy}} & \multicolumn{2}{c}{\textbf{Drake}} \\
\cmidrule(lr){2-3}\cmidrule(lr){4-5}\cmidrule(lr){6-7}
$N$ & time & residual & time & residual & time & residual \\
\midrule
$1$  & $1.57$\,s  & $0$                & $0.33$\,s  & $0$                & $0.44$\,s & $0$ \\
$2$  & $2.17$\,s  & $1\times10^{-16}$  & $0.38$\,s  & $4\times10^{-16}$  & $0.24$\,s & $2\times10^{-16}$ \\
$3$  & $21.63$\,s & $4\times10^{-16}$  & $0.59$\,s  & $8\times10^{-15}$  & $0.25$\,s & $6\times10^{-15}$ \\
$5$  & $7.45$\,s  & $2\times10^{-16}$  & $2.58$\,s  & $1\times10^{-15}$  & $0.24$\,s & $7\times10^{-14}$ \\
$8$  & $42.03$\,s & $2\times10^{-15}$  & \multicolumn{2}{c}{\emph{timeout}}  & $0.22$\,s & $5\times10^{-14}$ \\
$12$ & $64.43$\,s & $6\times10^{-15}$  & \multicolumn{2}{c}{---}             & $0.22$\,s & $1\times10^{-12}$ \\
$20$ & \multicolumn{2}{c}{\emph{timeout}}  & \multicolumn{2}{c}{---}        & $0.23$\,s & $1\times10^{-11}$ \\
$30$ & \multicolumn{2}{c}{---}             & \multicolumn{2}{c}{---}        & $0.23$\,s & $6\times10^{-11}$ \\
\bottomrule
\end{tabular}
\end{table}

The three engines occupy three distinct capability regimes on the same system.
The computer-algebra package returns equations to 5 links; the symbolic compiler
to 12; the recursive library to 30, the largest size tested. At 8 links one
engine returns nothing while the other two return correct equations; at 20, a
second returns nothing while the third continues.

Three observations follow.

The recursive library is in a different class of behaviour. Its cost is flat at
approximately $0.23$\,s from 5 links to 30, an $O(N)$ articulated-body algorithm
set against two symbolic approaches whose cost grows explosively. Its residual
grows with system size, from $6\times10^{-15}$ at 3 links to $6\times10^{-11}$
at 30, remaining four orders of magnitude inside the adjudication tolerance.

The symbolic compiler reaches further than the computer-algebra package driven
as documented --- 12 links against 5. This is a defensible comparative claim
against that engine used in the manner its documentation prescribes. It is not a
claim against the recursive library, which exceeds both by a wide margin.

Cost is not monotonic in system size for the symbolic compiler: three links cost
$21.63$\,s against $7.45$\,s for five. The effect is attributable to a threshold
in the engine's symbolic-to-numeric handling and is reported rather than
smoothed.

The finding that matters is what did not occur. Every engine that returned an
answer returned a correct one, and every engine that could not, failed by never
returning. A timeout is the loud failure mode: the researcher knows the
computation did not finish. The engines therefore differ in reach, not in
honesty, and the study's original hypothesis remains unsupported on this axis as
on every other.

% ===========================================================================
\section{A failure shared by every implementation, including the referee}
% ===========================================================================

\subsection{The configuration}
Consider the $N$-link chain with every joint angle set to $\pi$ and every
velocity set to zero: the chain stands inverted, each link balanced directly
above the one below. The exact solution is that nothing moves. Every term in the
equations of motion vanishes analytically. The gravity term is proportional to
$\sin\pi$, which is zero. The Coriolis term vanishes because the velocities are
zero and the angles equal. The configuration is an equilibrium, and the exact
trajectory is the constant function.

It is, of course, an \emph{unstable} equilibrium. That distinction is the whole
of what follows.

\subsection{The observation}
All three engines, and the independent reference, return approximately 20
radians of motion over a ten-second horizon. All four report success. None emits
a warning, a diagnostic, or any indication that the returned trajectory bears no
relation to the exact solution of the system as specified.

Energy is conserved throughout to $2.737\times10^{-10}$ --- the same tolerance
as in every other case in the study, and identical across all three engines
(Table~\ref{tab:amplitude}, final row). The trajectory is smooth, physically
plausible in isolation, and entirely wrong. \textbf{That energy is conserved so
well during the fall is precisely why every check in the suite passes it.}

\subsection{The mechanism}
The cause is arithmetic rather than algorithmic. In binary64, $\sin(\pi)$
evaluates not to zero but to $1.2246\times10^{-16}$, because $\pi$ is not
exactly representable. This leaves a residual gravitational acceleration of
$1.20\times10^{-15}$ at the initial state. The equilibrium being unstable, that
residual is amplified exponentially, and within the ten-second horizon it
reaches full scale.

No finite-precision computation can remain at an unstable equilibrium. The
observation is a property of the problem in floating-point arithmetic, not a
defect in any engine. That the independent reference --- which shares no code
with any engine and was derived by hand --- exhibits the identical behaviour is
what establishes this. Had the reference remained at rest, the finding would
have been a defect in three engines. It did not, and the finding is something
else.

\subsection{Why the case is retained}
An earlier version of this axis was written on the expectation that an engine
reporting motion at the inverted equilibrium is wrong and an engine reporting no
motion is right. That expectation was false. It is recorded here, and corrected
in the module rather than silently revised, because the ease with which it was
held is part of the finding: the intuition that a correct implementation should
return the exact answer at an exact equilibrium is natural, widely shared, and
incorrect. It is instance 4 of Table in \S\ref{sec:threats}.

What the case exposes is real and shared. Every implementation tested returns
success and hands back a large trajectory where the exact answer is no motion at
all, and none warns that the initial condition is an unstable equilibrium whose
evolution is dominated by round-off. This is the study's original thesis
appearing as a universal blind spot rather than as a difference between engines.
It is a weaker result than a disagreement would have been, and an honest one.

\subsection{What follows for differential testing}
Differential testing rests on the premise that independent implementations fail
independently, so that disagreement localises incorrectness. This case is a
counterexample of a specific kind. The failure is not independent across
implementations because its cause is not in any implementation: it is in the
interaction between the problem and the arithmetic every implementation shares.

Cross-implementation comparison cannot detect a failure of this kind, because
every implementation exhibits it identically. Adding a fourth engine, a fifth,
or a reference derived from first principles by an independent author would not
help --- this study added exactly such a reference, and it failed in the same
way.

The class of failures invisible to differential testing is therefore not empty,
and it is not exotic. Unstable equilibria are not pathological configurations
contrived for a test suite; they are the operating point of any
inverted-pendulum control problem and appear throughout robotics and
biomechanics. A method that cannot see them has a blind spot in a region
practitioners actually work in.

% ===========================================================================
\section{What these measurements do and do not support}
% ===========================================================================

\subsection{The integrator was pinned, which removes the variable}
Every engine supplied only its right-hand side, and integration was a single
call with identical tolerances for all of them. Energy drift is therefore
dominated by the shared integrator, not by the engines. Identical drift is
evidence that the equations are equivalent, which is what pinning was intended
to establish. The measurement is structurally incapable of ranking the engines
on integration accuracy, and that was its purpose.

\subsection{The differences are not differences}
The worst energy-drift values across the amplitude axis are $4.74$, $4.77$ and
$7.78\times10^{-10}$ --- the same quantity to within a factor of $1.6$. In the
chaotic case the three values differ by $0.2$ per cent relative, twelve orders
of magnitude below any physically meaningful quantity. No ranking can be built
on separations of this size.

\subsection{Divergence time cannot rank the engines}\label{sec:divergence}
In a chaotic system the time at which one trajectory departs from another is
governed by the size of the initial difference between their right-hand sides,
amplified exponentially. With the Lyapunov rate
$\lambda = 2.556\,\mathrm{s}^{-1}$ implied by the measured growth, a seed $s$
should cross a $10^{-6}$ threshold at
$t = \lambda^{-1}\ln(10^{-6}/s)$.

\begin{center}
\begin{tabular}{@{}lrrr@{}}
\toprule
\textbf{Engine} & \textbf{$N{=}3$ residual} & \textbf{Predicted} & \textbf{Observed} \\
\midrule
MechanicsDSL & $1.89\times10^{-15}$ & $7.86$\,s & $6.96$\,s \\
Drake        & $1.07\times10^{-14}$ & $7.18$\,s & $5.69$\,s \\
SymPy        & $2.62\times10^{-14}$ & $6.83$\,s & $6.28$\,s \\
\bottomrule
\end{tabular}
\end{center}

\noindent \textbf{The band is predicted; the ordering is not.} Predicted times
span $6.8$--$7.9$\,s and observed times span $5.7$--$7.0$\,s, so the mechanism
is right --- but the predicted order is MechanicsDSL, Drake, SymPy while the
observed order is MechanicsDSL, SymPy, Drake. Drake and SymPy transpose.

The reason the ordering fails is instructive, and it strengthens rather than
weakens the conclusion. The three seeds span barely one order of magnitude while
the amplification spans eleven. The effective seed is also the residual
\emph{along the actual trajectory}, not the maximum over random probe states,
and the threshold crossing is itself noisy in a chaotic system. Divergence time
therefore has no resolving power between engines whose equations agree this
closely: it measures the Lyapunov exponent, with the engine identity as a
rounding error.

\subsection{Why one engine's residual is larger, and why that is not worse}
The recursive library evaluates dynamics with an articulated-body algorithm
built for arbitrary topologies, contact, and real-time control. The reference
and the symbolic engines perform a direct solve of one hand-derived Lagrangian
for one fixed topology. More arithmetic operations accumulate more rounding. At
$7\times10^{-14}$ on the longer chains this is a deliberate
generality-for-precision trade, not a defect, and it remains six orders of
magnitude inside the adjudication tolerance. The cart--pole results of
Table~\ref{tab:cartpole}, where the same engine returns the smallest residuals
of the three, confirm the interpretation.

\subsection{Where the engines genuinely differ, the author's engine is behind}
\label{sec:behind}
Compile time at three links is approximately 37 times slower than the
computer-algebra package ($21.63$\,s against $0.59$\,s, Table~\ref{tab:scaling}).
The Hamiltonian pathway times out from three links, where the other engines do
not. On the near-singular axis its accuracy is $5.0\times10^{-12}$ against
$1.7\times10^{-15}$. On the cart--pole its residuals are two orders of magnitude
larger than the recursive library's. These measurements are reported for the
same reason as the rest: a study that reported only the comparisons its author's
engine wins would not be worth reading.

\subsection{What can be claimed}
That the three engines agree with an independent closed-form derivation, and
with one another, to within $\mathbf{3\times10^{-14}}$ on equations of motion,
and conserve energy indistinguishably under a pinned integrator, across
amplitudes spanning the regular and chaotic regimes, across three mechanism
families, and across constrained and unconstrained pathways.

The bound is set by the largest residual anywhere in the study --- the
computer-algebra package at $2.62\times10^{-14}$ on the three-link chain
(Table~\ref{tab:eom}) --- not by the smallest. It is stated as
$3\times10^{-14}$ rather than more tightly because a bound that does not cover
every measurement is not a bound.

That is an independently adjudicated correctness result. A claim of superiority
for any one engine would rest on differences of a fraction of a per cent in a
metric governed by a shared integrator, and --- in the case of the author's
engine --- would be advanced by the party with an interest in it. Such a claim
would be identified immediately and would discredit the work that is sound.

% ===========================================================================
\section{A diagnostic for convention mismatch}\label{sec:diagnostic}
% ===========================================================================

Three times during this study a comparison produced an error that was identical
across every case in a sweep, to many significant figures. On each occasion the
cause was a convention mismatch between the reference and the engine adapter,
not a defect in the engine.

\begin{center}
\begin{tabular}{@{}lL{0.36\textwidth}L{0.24\textwidth}@{}}
\toprule
\textbf{Observed error} & \textbf{Cause} & \textbf{Degenerate case that hid it} \\
\midrule
$4.53$, and exactly $1.00$ at dead centre & Hamiltonian pathway integrates $(q,p)$ and returns $\dot p$; compared against $\ddot q$ & $N{=}1$, where $M = I$ makes $p = \dot q$ \\
$2.5$--$6.4$, growing in $N$ & engine reports \emph{relative} joint angles; reference uses \emph{absolute} & $N{=}1$, where the two coincide \\
exactly $2.00$, all 20 cases & pole offset rotated about $+y$ mirrors the mechanism, flipping the cart's acceleration & --- \\
\bottomrule
\end{tabular}
\end{center}

The signature generalises. A discrepancy arising from genuine numerical
difference varies with configuration, because the quantities producing it vary
with configuration. A discrepancy that is constant across cases with different
parameters cannot have that origin. It indicates instead that two
implementations are being asked different questions, and that the comparison
rather than the implementation is at fault.

Under a relative-error metric the two commonest degenerate cases produce
characteristic values. A term absent entirely from one side yields a relative
error of exactly $1$; a term present with the wrong sign yields exactly $2$,
since $|a_d - a_r| = 2|a_r|$ divides out to exactly $2$ wherever
$|a_r| > 1$. Both were observed here.

A second tell is that the error vanishes in a degenerate configuration and grows
away from it. Genuine derivation errors rarely switch on at a particular system
size; convention mismatches characteristically do, because degenerate cases make
differing conventions coincide.

The rule is stated for use rather than for interest: \emph{in a
cross-implementation comparison, a discrepancy identical to many significant
figures across cases with different parameters indicates a convention mismatch
in the comparison, not a defect in the implementation under test.} Its practical
value is that it fires early, before a spurious defect is reported. Each of the
three instances here was resolved by inspecting the mass matrix directly, and
each would otherwise have been reported as a defect in widely used software.

% ===========================================================================
\section{Conclusion}
% ===========================================================================

Three independently developed rigid-body dynamics engines were compared under
adjudication by a closed-form reference sharing no library with any of them,
across a frozen adversarial case matrix, three structurally distinct mechanism
families, both dynamical regimes, and constrained and unconstrained pathways.

They agree to the limit of double-precision arithmetic. No engine returned an
incorrect answer while reporting success on any adjudicated case. Where they
differ, they differ in reach: the largest system each will answer varies by a
factor of six across the three, and every engine that failed did so by not
returning rather than by returning something wrong.

The single failure observed is shared by all three engines and by the
independent reference. At an unstable equilibrium every implementation reports
success while returning a trajectory the exact solution does not have, and none
warns. Its cause lies in the interaction between the problem and
finite-precision arithmetic, not in any implementation.

The consequence for method is the paper's main contribution. Differential
testing between implementations is a strong instrument for the failures
implementations make independently, and this study found none of those in three
mature engines. It is structurally blind to the failures they share. Where the
shared failure arises from arithmetic rather than from code, adding
implementations does not help, because the reference exhibits the failure too.
Detecting such cases requires reasoning about the problem --- here, recognising
that the initial condition is an unstable equilibrium --- rather than comparing
implementations of it.

That the engines agree is a reassuring result for practitioners who rely on
them. That they agree even when all of them are wrong is the more useful one.

% ===========================================================================
\section*{Disclosure}\label{sec:disclosure}
\addcontentsline{toc}{section}{Disclosure}
% ===========================================================================

The author maintains one of the three engines measured. This is disclosed in
every document of the study. The mitigation is structural: no engine adjudicates
itself or any other, and the referee shares no library with any of them. The
case matrix was frozen before measurement under a rule forbidding post hoc
removal of cases, and \S\ref{sec:behind} reports the measurements on which the
author's engine performs worst.

\paragraph{Use of AI assistance.}
An AI assistant (Anthropic Claude) was used for portions of the implementation
and analysis. The study design, the adjudication criteria, the freeze protocol,
and the manuscript are the author's. Every reported figure was produced by
executing committed code and checked against the run record.

\paragraph{Artefacts.}
Every figure reported was produced by executing committed code. Artefacts ---
the case matrix, the reference implementations, the engine adapters, the sweep
scripts, and the raw output records --- are available at
\href{https://doi.org/10.5281/zenodo.22315173}{\texttt{doi:10.5281/zenodo.22315173}},
archived under the MIT licence. The repository is at
\url{https://github.com/MechanicsDSL/mechanicsdsl} (study material under
\texttt{stress\_suite/}).

% ===========================================================================
\begin{thebibliography}{9}
\addcontentsline{toc}{section}{References}

\bibitem{barr2015}
E.~T. Barr, M.~Harman, P.~McMinn, M.~Shahbaz, and S.~Yoo.
\newblock The oracle problem in software testing: A survey.
\newblock \emph{IEEE Transactions on Software Engineering}, 41(5):507--525,
  2015. \newblock \doi{10.1109/TSE.2014.2372785}

\bibitem{erez2015}
T.~Erez, Y.~Tassa, and E.~Todorov.
\newblock Simulation tools for model-based robotics: Comparison of Bullet,
  Havok, MuJoCo, ODE and PhysX.
\newblock In \emph{IEEE International Conference on Robotics and Automation
  (ICRA)}, pages 4397--4404, 2015.

\bibitem{gonzalez2016}
F.~Gonz\'alez, D.~Dopico, R.~Pastorino, et~al.
\newblock Behaviour of augmented Lagrangian and Hamiltonian methods for
  multibody dynamics in the proximity of singular configurations.
\newblock \emph{Nonlinear Dynamics}, 85:1491--1508, 2016.
  \newblock \doi{10.1007/s11071-016-2774-5}

\bibitem{hatton1994}
L.~Hatton and A.~Roberts.
\newblock How accurate is scientific software?
\newblock \emph{IEEE Transactions on Software Engineering}, 20(10):785--797,
  1994.

\bibitem{iftomm}
IFToMM Technical Committee for Multibody Dynamics.
\newblock Library of computational benchmark problems.
\newblock \url{https://www.iftomm-multibody.org/benchmark}.
  \newblock Accessed September 2026.

\bibitem{kanewala2014}
U.~Kanewala and J.~M. Bieman.
\newblock Testing scientific software: A systematic literature review.
\newblock \emph{Information and Software Technology}, 56(10):1219--1232, 2014.
  \newblock \doi{10.1016/j.infsof.2014.05.006}

\bibitem{knight1986}
J.~C. Knight and N.~G. Leveson.
\newblock An experimental evaluation of the assumption of independence in
  multiversion programming.
\newblock \emph{IEEE Transactions on Software Engineering}, SE-12(1):96--109,
  1986.

\bibitem{mckeeman1998}
W.~M. McKeeman.
\newblock Differential testing for software.
\newblock \emph{Digital Technical Journal}, 10(1):100--107, 1998.

\end{thebibliography}

\end{document}
